\documentclass[10pt,letterpaper]{article}
\usepackage{cornell-notes}

\title{Clause 20.02.05: Pointer Alignment}
\author{}
\date{}
\setDocTitle{Clause 20.02.05: Pointer Alignment}
\setDocSubtitle{C++ 2024}
\setDocOwner{C++ 2024 Cornell Notes}
\setDocDate{\today}
\setCornellCollection{C++ 2024 Cornell Notes}
\setCornellUnitType{Clause}
\setCornellUnitNumber{20.02.05}
\setCornellUnitTitle{Pointer Alignment}
\hypersetup{pdftitle={Clause 20.02.05: Pointer Alignment},pdfsubject={Cornell notes},pdfauthor={}}

\begin{document}
\maketitle
\begin{CornellOverviewBox}{Overview}
\textbf{Classification:} Standard library - Normative\\[2pt]
\textbf{Learning objective:} Understand the rules, terminology, and semantic consequences associated with Pointer alignment.
\end{CornellOverviewBox}\begin{CornellSummaryBox}{Summary}
{\sffamily\bfseries\color{Accent} Summary}\par\vspace{3pt}Section 20.2.5 focuses on Pointer alignment. Apply the specified interface together with its mandates, effects, result conditions, exception behavior, and complexity constraints. When applying it, preserve the distinction between mandatory requirements, recommendations, permissions, and behavior deliberately left undefined, unspecified, or implementation-defined.
\end{CornellSummaryBox}

\end{document}
